\section{Regularity of the prescribed reached theta sections}
\label{sec:geometry-reached-regularity}

The missing regularity input in the primary trace is arithmetic.  A
finite-order operator need not preserve a specified integral lattice.
For the original theta operators, however, the degeneration gives an
exact coefficient valuation law.  That law proves regularity, including
at primes dividing the torsion order, once the original integral model
and its PEL uniformization have been supplied.

\subsection*{The precise supplier and its retained model}

Fix a complete discrete valuation ring $R$ of characteristic zero, with
fraction field $K$, uniformizer $\pi$, and the original integral framed
point.  Retain its smooth, separated, quasi-compact group model
$\mathcal G$, its rigidified line $\mathscr L_{\mathcal G}$, and its
PEL uniformization.  These are the actual objects in the source's
\texttt{Good3} contract, not replacements with a different integral
lattice.  In particular the uniformization retains the following data:
its abelian scheme $\mathcal B/R$, its cubical rigidified line
$\mathscr L$, its period lattice $X$ and map $c:X\to\mathcal B(R)$,
its positive monodromy quadratic form $Q:X\to\mathbb Z$, and the
coefficient-ring action with the original trace-one element $w_0$.
The original formal expansions on every $\pi^m$-thickening, their
unique generic-fibre counterparts, and their torsion identification
are included.  Constructing this uniformization is a separate remaining
geometric obligation.

For $\chi\in X$ and $k\geq1$, the actual coefficient line and lattice are
\begin{align*}
 \mathscr M_{k,\chi}
  &=\mathscr L^{\otimes(k-1)}\otimes t_{c(\chi)}^*\mathscr L
       \otimes(c(\chi)^*\mathscr L)^{-1}\otimes e^*\mathscr L,\\
 M_{k,\chi}&=H^0(\mathcal B,\mathscr M_{k,\chi}).
\end{align*}
The last two factors are pulled back from $\operatorname{Spec}R$.
Thus the four original exponents are $(k-1,1,-1,1)$.
Write $b(x,y)=Q(x+y)-Q(x)-Q(y)$.
The period identities give $b$ bilinear, $b(x,x)=2Q(x)$, and
$Q(x)>0$ for $x\ne0$.  The rank-zero case $X=0$ is allowed.

Extend $v(\pi)=1$ to an algebraic closure $\Omega$ of $K$.
For a coefficient $s_\chi$ over $\Omega$ write
$s_\chi\in M_{k,\chi}^{\geq a}$ if
$s_\chi=\sum_i z_i\sigma_i$, where $\sigma_i\in M_{k,\chi}$ and
$v(z_i)\geq a$.  All sums in this definition are finite.  For the
original Fourier family $s=(s_\chi)$ define
\begin{equation}\label{eq:reached-weighted-filtration}
 \mathcal N_k(s;r)\quad\Longleftrightarrow\quad
 s_\chi\in M_{k,\chi}^{\geq r+Q(\chi)/k}
                  \quad\hbox{for every }\chi\in X.
\end{equation}
This is exactly the source's \texttt{NuGe} filtration.  It uses the
integral coefficient modules; it is not a valuation on a chosen
rational coefficient basis.

\begin{theorem}\label{thm:reached-sections-regular}
Under the supplied model and uniformization above, every original
reached theta section is regular on $\mathcal G$.  Here ``reached''
means that it starts from the given framed level-one generator and is
obtained by finitely many of the following original operations:
normalized pullback by $a\in\mathcal O_E$ with
$a\bar a=N$, $N$ positive and odd, from level $k$ to level $kN$,
where $k=1$ or $k\geq4$; and canonical symmetric translation by a
level-$k$ torsion point.  The integral section restricts to the
prescribed generic section with scalar exactly one.

More precisely, the expansion of every such section satisfies
$\mathcal N_k(s;0)$.  No factor involving the residue characteristic,
the torsion order, the rank of $X$, or the coefficient-field degree is
introduced.
\end{theorem}

\subsection*{The exact theta valuation law}

We first derive the arithmetic input to translation.  A torsion lift in
the original uniformization is represented by an abelian point $m$ and
a period $y\in X$ with $[k]m=c(y)$, together with its biextension
multipliers.  Properness of $\mathcal B$ extends $m$ over the valuation
ring of $\Omega$.  Choose integral generators for the fibre lines in
these multipliers.  If $t_\chi$ is the scalar comparing the original
biextension multiplier at $\chi$ with its integral generator, then
\begin{equation}\label{eq:reached-biextension-ratio}
                   k\,v(t_\chi)=b(y,\chi).
\end{equation}
Here the valuation is the rational extension of $v(\pi)=1$, so no
rounding of $b(y,\chi)/k$ is permitted.

For completeness, this identity follows directly from the integral
cube law.  Take the $k$th tensor power of the multiplier for $(m,c(\chi))$
and use $[k]m=c(y)$.  The resulting period comparison is the product of
the original period generator with valuation
$kQ(y)+b(y,\chi)$, the inverse $k$th power of the period generator at
$y$ with valuation $-kQ(y)$, and the inverse $k$th power of the
rigidification, of valuation zero.  Every remaining factor is the
ratio of two integral generators of the same line, hence a unit of the
valuation ring.  Taking valuations gives
$k\,v(t_\chi)=kQ(y)+b(y,\chi)-kQ(y)$.
In particular this argument remains valid when the residue
characteristic divides $k$; it is an equality in $\mathbb Q$, not a
division inside $R$.

Let $w$ be the original scalar-line generator of the theta lift, and
write $w=c_w e_w$ with $e_w$ integral.  If the action is multiplied by
its original nonzero normalization scalar $\zeta$, the cubical
translation rule sends coefficient $\chi$ to coefficient $\chi+y$ and
shifts its scalar valuation by
\[
            \frac{b(y,\chi)}{k}+v(c_w)+v(\zeta).
\]
This assertion is two-sided.  Translation by $m$ and its inverse are
integral automorphisms after base change to the valuation ring, and
the integral cube-law generators are invertible.  Flat base change
for $H^0$ therefore transfers precisely the lattices used in
\eqref{eq:reached-weighted-filtration}.  To see that this uses the
original pullback and not a free choice of scalar, the cube-law
isomorphism has four line factors
$\mathscr L_{c(\chi)+m},\mathscr L_{c(\chi)},\mathscr L_m,
\mathscr L_0$ with exponents $(1,-1,-1,1)$; its multiplier is the
specified $t_\chi c_w$, and $\zeta$ is retained separately.
The inverse uses their exact inverses.

Completing the quadratic identity gives the original action $T$ the
following exact filtration law:
\begin{equation}\label{eq:reached-exact-translation-shift}
 \mathcal N_k(s;r)\ \Longleftrightarrow\
 \mathcal N_k(Ts;r+c_T),\qquad
 c_T=v(c_w)+v(\zeta)-\frac{Q(y)}k.
\end{equation}
Indeed
$Q(\chi+y)=Q(\chi)+Q(y)+b(y,\chi)$.
The comparison is uniform in $\chi$ and in the section $s$.

If $T^h=1$, iteration gives
$\mathcal N_k(s;r)\Longleftrightarrow\mathcal N_k(s;r+hc_T)$.
The set of rational $r$ for which $\mathcal N_k(s;r)$ holds is downward
closed.  A downward-closed subset of $\mathbb Q$ invariant under a
nonzero rational translation is either empty or all of $\mathbb Q$:
iterate in the positive direction, then use downward closure.
Consequently it is invariant under $c_T$, whether or not $c_T=0$.
Equation~\eqref{eq:reached-exact-translation-shift} thus implies
\[
                 \mathcal N_k(Ts;r)\Longleftrightarrow
                 \mathcal N_k(s;r).
\]
This proof includes the zero section and does not require attainment
of an infimum.  The original canonical symmetric theta lift has finite
order.  If one first chooses a different finite-order lift via the
torsion identification, the comparison scalar is a root of unity:
raising both lifts to a common order shows this directly.  Its
valuation is zero, so the same assertion holds for the prescribed
canonical move, with its original scalar.

\subsection*{The actual normalized pullback is integral}

Suppose $a\bar a=N$.  The coefficient at
$\beta=\bar a\chi$ of the original normalized pullback is
\begin{equation}\label{eq:reached-pullback-coefficient}
             s'_{\bar a\chi}=u_\chi\,a^*s_\chi;
             \qquad s'_\beta=0\quad(\beta\notin\bar aX).
\end{equation}
The map $\chi\mapsto\bar a\chi$ is injective since $X$ is torsion-free
and multiplication by $a$ composes with it to $[N]$.
The generator $u_\chi$ is integral with integral inverse.  Here is why
this is true for the prescribed normalization.

The integral trace-form identification with the original $w_0$ and
the Rosati identity construct the rigidified cubical isomorphism
$a^*\mathscr L\simeq\mathscr L^{\otimes N}$.
One can compute it on the Poincar\'e expression
$(\mathrm{id},\lambda w_0)^*\mathscr P$:
$\bar a w_0a=Nw_0$ gives the isomorphism by bilinearity.
Equivalently the integral cube law and the trace-form identification
give the same isomorphism on the original line.  Applying this to
$\mathscr M_{k,\chi}$, and using
$c(\bar a\chi)=\bar a\,c(\chi)$ and the Rosati biextension identity,
constructs an integral generator
\[
 e_\chi\in
 H^0\!\left(\mathcal B,
   \mathscr M_{kN,\bar a\chi}\otimes(a^*\mathscr M_{k,\chi})^{-1}
       \right)^\times.
\]
Its value at the zero section is the given rigidification to the
power $k(N-1)$.  The original generic $u_\chi$ has exactly this value.
Their ratio is a global invertible function on the proper geometrically
connected abelian variety $\mathcal B_\Omega$, hence a scalar; its value
at zero is one.  Therefore $u_\chi=e_\chi$ exactly.
There is no extra integral scalar or denominator to estimate.

The balanced period pairing gives
$b(\bar a\chi,\bar a\chi)=b(\chi,a\bar a\chi)=Nb(\chi,\chi)$.
Since $b(x,x)=2Q(x)$, equality of integers gives
\[
                 Q(\bar a\chi)=NQ(\chi).
\]
Thus the source threshold in \eqref{eq:reached-pullback-coefficient}
is unchanged:
$Q(\bar a\chi)/(kN)=Q(\chi)/k$.
Integral pullback and the integral generator prove
\begin{equation}\label{eq:reached-pullback-filtration}
       \mathcal N_k(s;r)\ \Longrightarrow\
       \mathcal N_{kN}(s';r).
\end{equation}
This is a forward statement, exactly what reaching a section requires;
no inverse pullback or index bound is assumed.

\subsection*{The initial coefficient and induction}

The supplied framed level-one section is integral.  Its coefficient at
$0$ lies in $M_{1,0}$: the integral formal expansion has precisely this
coefficient, and the original line identification and framed section
agree with scalar one.  The period recurrence moves a coefficient at
$\alpha$ to that at $\alpha+kx$ by an integral line isomorphism and a
scalar of valuation
\[
                         kQ(x)+b(x,\alpha).
\]
This is the same period comparison used above.  It agrees with the
quadratic threshold because
\[
 Q(\alpha+kx)-Q(\alpha)=k^2Q(x)+kb(x,\alpha).
\]
For $k=1$, every element of $X$ is obtained from $0$.  Since $Q(0)=0$,
the integral zero coefficient proves $\mathcal N_1(s;0)$ for the
original framed generator.  Applying
\eqref{eq:reached-pullback-filtration} and finite-order translation
invariance inductively proves $\mathcal N_k(s;0)$ for every reached
section.  All intermediate levels and the source restriction
$k=1$ or $k\geq4$ are retained.

\subsection*{Coefficient descent and integral extension}

The remaining issue is genuine integral descent, not rational
uniqueness.  We give both steps explicitly.

First suppose a coefficient defined over $K$ is a finite sum
$\sum_i z_i\sigma_i$ with $v(z_i)\geq0$ and
$\sigma_i\in M_{k,\chi}$.  The finite $R$-module
$V=R+\sum_iRz_i\subset\Omega$ is torsion-free, hence free.
The element $1$ is primitive in $V$: otherwise
$1=\pi z$ with $z\in V$, contradicting $v(z)\geq0$.
Extend $1$ to an $R$-basis and take its coordinate projection
$\ell:V\to R$.  Extend $\ell$ $K$-linearly to $\Omega$; it fixes $K$.
Applying it to the coefficient identity gives
$s_\chi=\sum_i\ell(z_i)\sigma_i\in M_{k,\chi}$.
This works for an arbitrary valued extension and needs neither an
integral trace nor division by its degree.

Because $Q(\chi)\geq0$, the already proved $\mathcal N_k(s;0)$ makes
every coefficient integral over the valuation ring.  The preceding
argument puts each original $K$-coefficient in its exact $R$-module
$M_{k,\chi}$.

Now let $s$ be the original section over $K$.  Quasi-compactness and
separatedness give an $m\geq0$ such that $\sigma=\pi^m s$ extends to
$\mathcal G$: on a finite affine cover one clears its denominators.
The integral expansion of $\sigma$ has every coefficient divisible by
$\pi^m$.  If $m>0$, reduce the original expansion formula on the
$\pi^m$-thickening.  Every coefficient is zero, so that formula makes
$\sigma$ itself zero on this thickening.  The sheaf exact sequence
\[
 0\longrightarrow\mathscr L_{\mathcal G}^{\otimes k}
   \xrightarrow{\ \pi^m\ }
   \mathscr L_{\mathcal G}^{\otimes k}
   \longrightarrow
   \mathscr L_{\mathcal G}^{\otimes k}/\pi^m
   \longrightarrow0
\]
is exact because the smooth model is flat.  Left exactness of global
sections supplies $\tau$ with $\sigma=\pi^m\tau$.
On the generic fibre, cancellation gives $\tau=s$.  For $m=0$ no
cancellation is necessary.  This proves Theorem~\ref{thm:reached-sections-regular}
with scalar one.  The argument uses the reconstruction formula on
each integral thickening; rational Fourier uniqueness alone would
not supply this divisibility.

\subsection*{Effect on the primary trace and remaining frontier}

If the tuple $F$ in Theorem~\ref{thm:full-primary-theta-trace} is reached,
then each of its prescribed translates $G_u$ is reached by definition.
Theorem~\ref{thm:reached-sections-regular} supplies the exact integral
sections needed there.  Their product invariance under the stabilizer
is the section identity proved before formal trace, so no
harmonic-value-only invariance is substituted.  Normalized pullback
in the adapted lattice argument again produces reached sections.
Accordingly the primary theta denominator estimate now needs no
separate reached-section regularity assumption once the original
\texttt{Good3} model is supplied.

This closes the reached-section regularity subobligation under that
model contract.  It does not construct the global integral model or
its degeneration, prove the mixed level-torsion identities needed for
the other local branch, or supply theta spanning and the raw global
norm estimate.  The conductor-one specialization does not remove
these retained geometric inputs.  Exact source interfaces and bounded
excerpt hashes are recorded in the accompanying certificate
\cite{reached-regularity-source}.
